%% kapa: exam=1 round=29 year=2018 period=1 form=A subject=economics subject_ko=경제학원론 seq=4 number=44
%% source: Q-Net 감정평가사 기출문제 (cst00309) artlSeq=5202808 file=extracted/1008/5202808/2018년 제29회 감정평가사 1차시험 기출문제/2018년 제29회 감정평가사 1차시험 기출문제/1교시 A형.pdf
%% conversion: kordoc 4.24.51 -> tools/gen1.py (auto); review=proofread-sonnet; flags=none
\begin{question}{44}
주유소에서 휘발유를 구입하는 모든 소비자들은 항상 “5만 원어치 넣어주세요”라고 하는 반면, 경유를 구입하는 모든 소비자들은 항상 “40리터 넣어주세요”라고 한다. 현재의 균형상태에서 휘발유의 공급은 감소하고, 경유의 공급이 증가한다면, 휘발유 시장과 경유 시장에 나타나는 균형가격의 변화는? (단, 휘발유 시장과 경유 시장은 완전경쟁시장이며, 각 시장의 공급곡선은 우상향하고, 다른 조건은 일정하다.)
\begin{choices}
  \item 휘발유 시장: 상승, 경유 시장: 상승
  \item 휘발유 시장: 상승, 경유 시장: 하락
  \item 휘발유 시장: 하락, 경유 시장: 불변
  \item 휘발유 시장: 하락, 경유 시장: 하락
  \item 휘발유 시장: 불변, 경유 시장: 불변
\end{choices}
\end{question}
\answer{2}
